\makeatletter
\def\input@path{{./}{../}{../../}{../../../}{../../../../}{preamble/}{../preamble/}{../../preamble/}{../../../preamble/}}
\makeatother
\input{preamble_exams.tex}
\input{graphs/highlights.tex}
% \GTenHighlightsfalse % Uncomment to disable instructional highlights.
\input{graphs/specific_highlights.tex}

\hypersetup{pdftitle={Grade 10 Global Economics: Term 3, Lesson 1, C7},pdfsubject={Home practice: Inverse Normal Reasoning and Value at Risk}}
\title{Grade 10 Global Economics\\Term 3, Lesson 1: Home Practice\\C7 --- Inverse Normal Reasoning and Value at Risk}
\author{}
\date{}
\begin{document}
\maketitle
\setcounter{tocdepth}{1}
\setcounter{secnumdepth}{2}
\phantomsection\label{toc}
\tableofcontents
\section*{Evaluated criteria}
\begin{description}
\item[\textbf{C7}] Applies inverse normal reasoning to find values from given probabilities, interpreting percentiles, confidence-related values, and value at risk in financial contexts.
\end{description}
\textbf{How to practise.} Write the left-cumulative target before choosing a z-value. State whether the answer is an observation cutoff, an interval of observations, a critical multiplier or a loss threshold. Try each task on paper before opening its solution link. Record the model, working, units and an interpretation; use the arrows to return.
All datasets and economic or financial models in this activity are illustrative, rather than observed market evidence.

\clearpage
\section{Introduction}\label{sec:introduction}
\subsection*{Highlight legend}
\phantomsection\label{tab:highlight-legend}
\begin{tabularx}{\linewidth}{@{}>{\raggedright\arraybackslash}p{.20\linewidth}>{\raggedright\arraybackslash}p{.22\linewidth}>{\raggedright\arraybackslash}X@{}}
\toprule
\textbf{Symbol or notation} & \textbf{Concept} & \textbf{Brief meaning}\\
\midrule
$\HLCumulativeProbability{p},\ \HLCriticalValue{z_p}$ & \HLCumulativeProbability{Cumulative target} and \HLCriticalValue{cutoff} & Area and its standard position. \\
$\GTenHLValue{x_p}=\mu+\sigma z_p$ & \GTenHLValue{Observation percentile} & Cutoff in original units. \\
$\HLConfidence{C},\ \HLCriticalValue{z^*}$ & \HLConfidence{Central coverage} and \HLCriticalValue{multiplier} & Middle area and its boundary. \\
$\mathrm{SE}$ & Standard error & Spread of an estimator. \\
$R,\ L,\ \mathrm{VaR}_{C}$ & Return, loss and value at risk & A horizon-specific loss quantile. \\
\bottomrule
\end{tabularx}
\subsection{Introductory example 1: A high spending percentile}\label{ex:intro-one}
For illustrative weekly spending $X\sim N(100,20^2)$ dollars, the \HLCumulativeProbability{95th percentile} has left area $\HLCumulativeProbability{0.95}$. The standard cutoff is $\HLCriticalValue{z_{0.95}\approx1.645}$, so $x=100+20(1.645)=132.90$ dollars. This is a cutoff for individual weekly observations.
\subsection{Introductory example 2: A confidence-related multiplier}\label{ex:intro-two}
A supplied estimator model states that $Z=(\widehat\mu-\mu)/\mathrm{SE}\sim N(0,1)$, where $\mathrm{SE}$ describes the spread of an estimator across repeated samples. A \HLConfidence{central coverage} of $\HLConfidence{95\%}$ leaves $\HLTailArea{2.5\%}$ in each tail, giving critical multipliers $\HLCriticalValue{\pm1.96}$. Here the scale is the estimator's standard error, rather than the SD of individual spending observations.
\section{Notation Reference and Practice Procedure}\label{sec:notation-and-procedure}
\subsection{Notation reference}\label{tab:notation-reference}
\begin{tabularx}{\linewidth}{@{}>{\raggedright\arraybackslash}p{.16\linewidth}>{\raggedright\arraybackslash}p{.18\linewidth}>{\raggedright\arraybackslash}p{.27\linewidth}>{\raggedright\arraybackslash}X@{}}
\toprule
\textbf{Symbol or notation} & \textbf{Concept} & \textbf{Definition} & \textbf{Meaning in the introductory examples}\\
\midrule
$\HLCumulativeProbability{p},\ \HLCriticalValue{z_p}$ & \HLCumulativeProbability{Cumulative target} and \HLCriticalValue{cutoff} & $Phi(z_p)=p$; an upper-tail probability $q$ uses $p=1-q$. & \textbullet\ \hyperref[ex:intro-one]{Example 1}: The individual 95th percentile uses $p=0.95$. \\
$\GTenHLValue{x_p}=\mu+\sigma z_p$ & \GTenHLValue{Observation percentile} & Value below which fraction $p$ of the model's observations lie. & \textbullet\ \hyperref[ex:intro-one]{Example 1}: The spending cutoff is 132.90 dollars. \\
$\HLConfidence{C},\ \HLCriticalValue{z^*}$ & \HLConfidence{Central coverage} and \HLCriticalValue{multiplier} & For symmetric central area $C$, $z^*=z_{(1+C)/2}$. & \textbullet\ \hyperref[ex:intro-two]{Example 2}: 95 percent central coverage needs a 0.975 upper cumulative target. \\
$\mathrm{SE}$ & Standard error & SD of the estimator across repeated samples; supplied when needed here. & \textbullet\ \hyperref[ex:intro-two]{Example 2}: It scales a confidence-related margin for a mean. \\
$R,\ L,\ \mathrm{VaR}_{C}$ & Return, loss and value at risk & For return in percentage points, $L=-VR/100$; $mathrm{VaR}_C$ is the $C$ quantile of loss. & \textbullet\ \hyperref[ex:intro-one]{Example 1}: A spending percentile is an observation cutoff; VaR likewise cuts off a distribution, but for financial loss. \\
\bottomrule
\end{tabularx}
\subsection{Practice procedure}
\begin{enumerate}
\item Identify the variable, units, normal mean and SD, and whether the target describes an individual observation or an estimator.
\item Convert an upper-tail statement to $p=1-q$. For a central area $C$, use lower and upper targets $(1-C)/2$ and $(1+C)/2$.
\item Read the supplied inverse values and transform with $x=\mu+\sigma z$. For a supplied estimator model use its SE in place of the individual-observation SD.
\item For VaR, state horizon, confidence, portfolio value and loss convention. Convert the lower return quantile into a positive loss amount and explain the remaining exceedance probability.
\end{enumerate}
\subsection{Inverse-normal reference values}
$\Phi(z)=P(Z\le z)$. The following approximate values supply every inverse lookup needed here:
\[
\begin{array}{c|rrrrrrrrr}
p&.01&.025&.05&.10&.90&.95&.975&.99&.995\\ \hline
z_p&-2.326&-1.960&-1.645&-1.282&1.282&1.645&1.960&2.326&2.576
\end{array}
\]
Use $N(\mu,\sigma^2)$ with variance as the second parameter. A central interval for observations describes individual outcomes. A confidence procedure for a mean instead uses an estimator's sampling distribution; the normal estimator model and SE are provided, so no new estimation method is required.
\section{General Worked Problems}\label{sec:general-worked}
\SectionProblemsTableOfContents{
\SectionProblemLink{prob:sec:general-worked0}{The top five percent of spending}\\
\SectionProblemLink{prob:sec:general-worked1}{A central interval changes with its coverage}\\
\SectionProblemLink{prob:sec:general-worked2}{A confidence-related critical value on an estimator scale}}

\subsection{The top five percent of spending}\label{prob:sec:general-worked0}
Weekly spending follows $X\sim N(100,20^2)$ dollars. A budget adviser wants the cutoff for the highest $5\%$ of weekly observations and the cutoff for the lowest $10\%$.
\QuestionsTasks[sec:general-worked]{{7}/{sol:sec:general-worked0}/{State each cumulative target and find both cutoffs. Interpret them without treating either as a confidence interval for a mean.}}
\subsubsection*{Worked solution}\label{sol:sec:general-worked0}
For the top $5\%$, $P(X>x)=0.05$, hence $P(X\le x)=0.95$ and $\HLCriticalValue{z=1.645}$. Invert standardisation:
\[
(x-100)/20=1.645,\quad x=100+20(1.645)=\GTenHLValue{132.90}\text{ dollars}.
\]
For the bottom $10\%$, the cumulative target is $0.10$, so
\[
x=100+20(-1.282)=\GTenHLValue{74.36}\text{ dollars}.
\]
About $5\%$ of individual household-weeks exceed 132.90 dollars, and $10\%$ are below 74.36 dollars. Neither cutoff estimates uncertainty about the population mean.
\secInicio[prob:sec:general-worked0][sec:general-worked]


\subsection{A central interval changes with its coverage}\label{prob:sec:general-worked1}
An annual investment return $R\sim N(6,2^2)$ is measured in percentage points; 6 means a return of $6\%$ of initial capital.
\QuestionsTasks[sec:general-worked]{{7}/{sol:sec:general-worked1}/{Find the central 90 percent and central 95 percent intervals for annual returns. Explain why increasing coverage changes both endpoints.}}
\subsubsection*{Worked solution}\label{sol:sec:general-worked1}
For central $\HLConfidence{90\%}$, each tail has $(1-0.90)/2=0.05$, so the standard endpoints are $\pm1.645$:
\[
[6-2(1.645),6+2(1.645)]=[2.71,9.29]
\]
in percentage points. For central $\HLConfidence{95\%}$, each tail is $0.025$, so use $\pm1.96$:
\[
[6-2(1.96),6+2(1.96)]=[2.08,9.92].
\]
The greater coverage requires more area on both sides of the centre and therefore a wider interval. These are intervals of modeled annual returns, not confidence intervals estimating the mean return. Approximately $5\%$ of observations lie outside the second interval, split equally between its tails.
\secInicio[prob:sec:general-worked1][sec:general-worked]


\subsection{A confidence-related critical value on an estimator scale}\label{prob:sec:general-worked2}
A supplier study reports an estimated mean invoice amount $\widehat\mu=50$ dollars with a supplied standard error of 2 dollars. Assume the estimator pivot $Z=(\widehat\mu-\mu)/2$ follows $N(0,1)$. A central 95 percent confidence procedure is specified as $\widehat\mu\pm z^*\mathrm{SE}$.
\QuestionsTasks[sec:general-worked]{{7}/{sol:sec:general-worked2}/{Choose the correct critical value and calculate the interval. Explain why 1.645 and the standard deviation of individual invoices would answer different questions.}}
\subsubsection*{Worked solution}\label{sol:sec:general-worked2}
Central $95\%$ leaves $2.5\%$ in each tail. Thus the upper cumulative target is $0.975$, giving $\HLCriticalValue{z^*=1.96}$. The margin is $1.96(2)=3.92$ dollars, so
\[
[50-3.92,50+3.92]=[46.08,53.92]\text{ dollars}.
\]
Under the supplied sampling assumptions, $95\%$ of intervals constructed this way in repeated studies contain the fixed mean invoice amount. This is different from saying that $95\%$ of individual invoices lie here. The value $1.645$ corresponds to a one-sided $95\%$ boundary or a central $90\%$ band. The individual-invoice SD concerns observation spread; this procedure specifically supplies the estimator SE.
\secInicio[prob:sec:general-worked2][sec:general-worked]
\section{Advanced Problems}\label{sec:advanced}
\SectionProblemsTableOfContents{
\SectionProblemLink{prob:sec:advanced0}{One-day value at risk from a return percentile}\\
\SectionProblemLink{prob:sec:advanced1}{Keeping the three interpretations separate}}

\subsection{One-day value at risk from a return percentile}\label{prob:sec:advanced0}
A portfolio starts at $V=10{,}000$ dollars. Its \textbf{one-day} return $R\sim N(0.10,1.00^2)$ is in percentage points: $R=1$ means a $1\%$ gain. Define monetary loss as $L=-VR/100$, so a negative return gives a positive loss. Define $\mathrm{VaR}_{95}$ as the loss threshold with $P(L\le\mathrm{VaR}_{95})=0.95$.
\QuestionsTasks[sec:advanced]{{7}/{sol:sec:advanced0}/{Find the one-day 95 percent VaR and then the one-day 99 percent VaR. Explain the tail reversal and why VaR is not a maximum possible loss.}}
\subsubsection*{Worked solution}\label{sol:sec:advanced0}
Large losses correspond to low returns because $L=-VR/100$. Thus the $95$th loss percentile comes from the $5$th return percentile:
\[
r_{0.05}=0.10+1.00(-1.645)=-1.545
\]
percentage points, a return of $-1.545\%$. Converting this return to loss gives
\[
\mathrm{VaR}_{95}=-10{,}000(-1.545)/100=\GTenHLValue{154.50}\text{ dollars}.
\]
Under the model, $P(L>154.50)\approx0.05$ over one day. For confidence $99\%$, use the $1$st return percentile:
\[
r_{0.01}=0.10-2.326=-2.226,\quad
\mathrm{VaR}_{99}=-10{,}000(-2.226)/100=\GTenHLValue{222.60}\text{ dollars}.
\]
Higher confidence moves deeper into the loss tail. Neither threshold is a maximum: the model still permits larger losses, with exceedance probabilities approximately $5\%$ and $1\%$. These daily numbers do not automatically apply to a week or year, and they depend on the assumed return distribution.
\secInicio[prob:sec:advanced0][sec:advanced]


\subsection{Keeping the three interpretations separate}\label{prob:sec:advanced1}
A report includes these statements: (A) an individual's 95th percentile uses $z=1.645$; (B) a central 95 percent observation interval uses $z=\pm1.96$; (C) a normal confidence procedure for a mean uses $\widehat\mu\pm1.96\,\mathrm{SE}$.
\QuestionsTasks[sec:advanced]{{7}/{sol:sec:advanced1}/{Explain why all three may be correct and identify the probability distribution and scale behind each statement.}}
\subsubsection*{Worked solution}\label{sol:sec:advanced1}
A sets left area $0.95$ for a single observation and transforms with its model mean and SD. B puts $0.025$ in each tail of the observation distribution, then transforms both standard cutoffs with the observation SD. C uses the same central standard-normal area as B, but its distribution is the estimator's repeated-sample distribution and its scale is SE. The numbers $1.645$ and $1.96$ answer different area questions; identical $1.96$ multipliers do not make observation intervals and confidence intervals interchangeable. A VaR statement additionally requires a loss convention, horizon and confidence level.
\secInicio[prob:sec:advanced1][sec:advanced]

\end{document}

